% ============================================================
% CONCLUSION
% ============================================================
\chapter*{Conclusion}
\pagestyle{styleGlobal}
\label{sec:conclusion}
\markboth{Conclusion}{Conclusion}
\addcontentsline{toc}{chapter}{Conclusion}

Au terme de ce travail, nous avons conçu et réalisé «~SALISA~», une application cross-plateforme de mise en relation entre prestataires et clients pour services professionnels, développée pour le compte d'Akieni.

Ce travail nous a permis de parcourir l'ensemble du cycle de développement logiciel, depuis l'analyse des besoins jusqu'à la réalisation d'une solution fonctionnelle, en passant par la conception architecturale et la modélisation UML selon la démarche \textbf{2TUP}. Nous avons pu mesurer concrètement la complémentarité entre un langage de modélisation expressif comme UML et un processus de développement structuré comme le 2TUP, dont l'approche en deux branches parallèles s'est révélée particulièrement adaptée à la complexité de «~SALISA~».

Sur le plan fonctionnel, «~SALISA~» propose un ensemble de fonctionnalités répondant aux besoins réels du marché congolais : mise en correspondance intelligente entre prestataires et demandeurs, système de confiance à plusieurs niveaux, paiement sécurisé par Mobile Money via un mécanisme d'escrow, messagerie en temps réel et notifications multi-canaux. L'application est conçue comme une \textit{Progressive Web Application} (PWA), optimisée pour les réseaux mobiles et accessible sans installation.

Ce travail illustre la capacité du numérique à structurer des marchés informels et à créer de la valeur économique locale, tout en répondant aux contraintes spécifiques du contexte africain. Il constitue une contribution concrète à la transformation numérique en cours en République du Congo.

En perspective, plusieurs axes d'amélioration pourront être envisagés : le déploiement d'un module de machine learning pour affiner l'algorithme de mise en correspondance, l'extension de «~SALISA~» vers d'autres villes du pays comme Pointe-Noire, le développement d'une application mobile native, et la mise en place d'un programme d'ambassadeurs pour accélérer son adoption.

Ce travail, mené dans le cadre de notre formation en Licence Professionnelle Génie Logiciel à l'ESGAE, nous a permis de consolider nos compétences techniques et méthodologiques, tout en relevant un défi concret ancré dans les réalités de notre pays.

\vspace{1cm}
\begin{flushright}
    \textit{Ko salisa biso nyonso} - Nous nous entraidons tous.
\end{flushright}
